\section{Benchmark and Experimental Setup}

\paragraph{P0 adapted retrieval-memory comparator.}
The P0 evaluation adapts an AgentPoison-style full-ReAct retrieval-memory setting. It measures raw poisoned retrieval, exposed poisoned retrieval, attack manifestation, utility, and defense interventions. This is an adapted comparator, not an official AgentPoison reproduction.

\paragraph{P1 deterministic synthetic MAS.}
The deterministic P1 benchmark uses a synthetic enterprise-assistant task and controlled multi-agent topologies. Provider calls are disabled. The strengthened deterministic matrix is used for topology and indirect-attack mechanism claims.

\paragraph{MiniMax-backed synthetic-runtime coverage.}
The canonical MiniMax-backed multi-agent synthetic-runtime coverage experiment uses three topologies: \texttt{chain\_4}, \texttt{star\_4}, and \texttt{blackboard\_4}. It uses seven attack settings: no attack, summary poisoning direct/indirect, workspace poisoning direct/indirect, and communication hijack direct/indirect. It compares four defenses: no defense, static ACL, prompt filter, and FlowFence-Lite. The matrix contains $3 \times 7 \times 4 \times 3 = 252$ configured runs with MiniMax final-writer calls.

\paragraph{Non-oracle held-out validation.}
The non-oracle validation adds held-out paraphrase attacks and evaluates \texttt{flowfence\_lite\_nonoracle}. The deterministic matrix contains 540 runs with provider calls disabled. The targeted MiniMax-backed synthetic-runtime validation contains 72 runs over held-out paraphrases, two topologies, four defenses, and three seeds.

\paragraph{Evidence handling.}
Raw traces, provider outputs, prompts, event JSONL, policy JSONL, and per-run metrics are not included in the paper draft. The draft uses committed high-level summaries, reviewed paper-facing tables, claims, and evidence-boundary documents.
